\chapter{РЕЗУЛЬТАТЫ И ОБСУЖДЕНИЕ}

В настоящей главе представлены условия проведения экспериментов,
результаты предобучения адаптивного токенизатора в составе языковой
модели на референсном геноме человека, результаты оценки на наборах
прикладных задач, а также анализ выученной токенизации в сопоставлении
с биологическими аннотациями. На момент сдачи работы полный
экспериментальный цикл находится в стадии выполнения; в разделах с
численными результатами фиксируется формат отчетности и протокол их
получения, а итоговые значения подлежат заполнению по факту завершения
обучения.

\section{Условия эксперимента}

В качестве основного предмета экспериментальной оценки выступает
адаптивный токенизатор в составе модели, описанной во второй главе:
размерность представлений токенов $d_{\text{tok}} = 64$, размерность
бэкбона $d = 256$, конфигурация Transformer~--- 6 слоёв, 8 голов,
размерность скрытого слоя FFN $4d = 1024$, dropout $0{,}1$. Конвейер
слияний адаптивно подбирает число уровней рекурсии под длину входа:
для $L = 2048$ при целевом $L_{\text{target}} = 1024$ выполняется один
уровень слияния, при котором выбирается до $30\%$ соседних пар.
Декодер на нуклеотидном уровне~--- per-bp с тремя вкладами (контекст
спана, эмбеддинг запросной позиции, skip-связь от исходных вложений) и
LayerNorm на сумме. Полное число обучаемых параметров составляет
порядка 5\,M (бэкбон ${\sim}4{,}85$\,M, токенизатор ${\sim}64$\,тыс.,
декодер ${\sim}87$\,тыс.).

В качестве обучающего корпуса был использован референсный геном
человека T2T-CHM13
v2.0~\cite{nurgalievT2TCHm13CompleteHuman2022} без Y-хромосомы. Сплит
выполнен по протоколу
HyenaDNA~\cite{nguenHyenaDNALongRangeGenomic2024}: chr8~--- тест,
chr9~--- валидация, остальные аутосомы и X-хромосома~--- обучение.
Длина обучающего окна $L = 2048$~п.н., шаг 2048 (без перекрытия);
фрагменты с долей символа $N$ выше 50\% пропускаются на уровне
семплера. Аугментация обратным комплементом активируется через флаг
конфигурации с вероятностью $0{,}5$. Размер пакета 64; накопление
градиентов отсутствует.

Оптимизация осуществляется при помощи
AdamW~\cite{loshchilovDecoupledWeightDecay2019} ($\beta_1 = 0{,}9$,
$\beta_2 = 0{,}95$, weight decay $0{,}1$); пиковая скорость
$3 \times 10^{-4}$ достигается за 200 шагов прогрева, после чего по
косинусной кривой снижается до $0{,}1$ от пика. Норма градиента
ограничена $1{,}0$. Смешанная точность обучения~--- \texttt{bfloat16}
через \texttt{torch.amp.autocast} (без \texttt{GradScaler}). Маскирование
по протоколу BERT~\cite{devlinBERTPretrainingDeep2019}: 15\% позиций
маскируются по схеме 80--10--10. Полный цикл предобучения~--- $10^4$
итераций.

Все запуски выполнены на одном GPU NVIDIA с поддержкой
\texttt{bfloat16} (архитектура Ampere или Hopper, не менее 24 ГБ
памяти). Логирование метрик и артефактов~--- через
MLflow~\cite{Zaharia_Accelerating_the_Machine_2018} с локальным
sqlite-бэкендом; конфигурация запуска и итоговый Hydra-конфиг
сохраняются вместе с чекпойнтами под общим корнем \texttt{outputs/}.

\section{Результаты предобучения}

Контроль хода обучения осуществляется по следующему набору метрик:
кросс-энтропия MLM на маскированных позициях, перплексия
$\mathrm{PPL} = \exp(\mathcal L_{\text{MLM}})$, бит на нуклеотид
$\mathrm{BPN} = \mathcal L_{\text{MLM}} / \ln 2$, норма градиента,
текущая скорость обучения, доля маскированных позиций в батче, а
также пропускная способность в нуклеотидах в секунду. На валидационной
хромосоме chr9 с интервалом 1000 шагов вычисляется \texttt{val/loss};
улучшение этой метрики является сигналом для сохранения чекпойнта
\texttt{best}.

Стоит отметить ожидаемое поведение метрик в начале обучения. На
первых нескольких сотнях шагов прогрева MLM-перплексия снижается с
начального значения $|V| = 9$ до уровня посимвольной модели; начиная
со 2--3 тысячи шагов, эффект конвейера слияний должен проявляться в
двух признаках~--- появлении неравномерного распределения длин спанов
$\ell_j$ и в снижении $\mathrm{BPN}$ относительно посимвольной базовой
линии при том же числе параметров. В случае коллапса конвейера
слияний (все спаны фактически имеют одинаковую длину) данные
признаки наблюдаться не будут, что является диагностическим
маркером необходимости пересмотра стратегии обучения.

Сравнение с базовыми моделями приведено в таблице~\ref{tab:pretrain}.
Все базовые модели обучаются с теми же гиперпараметрами оптимизатора
и при сопоставимом числе параметров.

\begin{table}[h]
\centering
\caption{Метрики предобучения на CHM13 (валидация на chr9, тест на chr8)}
\label{tab:pretrain}
\begin{tabular}{lcccc}
\toprule
Модель & PPL & BPN & TR & CR \\
\midrule
Char-Transformer & TBD & TBD & 1{,}00 & 0{,}00 \\
BPE-Transformer (4096) & TBD & TBD & TBD & TBD \\
Fixed $k$-mer ($k = 6$) & TBD & TBD & 0{,}167 & 0{,}833 \\
Char-Mamba & TBD & TBD & 1{,}00 & 0{,}00 \\
MxDNA~\cite{qiaoMxDNAAdaptiveDNA2024} & TBD & TBD & TBD & TBD \\
Адаптивный (предлагаемый) & TBD & TBD & TBD & TBD \\
\bottomrule
\end{tabular}
\end{table}

Прикладной критерий приемлемости был зафиксирован еще до запуска
экспериментов: $\mathrm{TR}_{\text{adapt}} - \mathrm{TR}_{\text{BPE}}
\leq 0{,}05$ и $\Delta \mathrm{BPN} \leq 0{,}02$. Предпочтительный
режим~--- $\mathrm{TR}_{\text{adapt}} \leq \mathrm{TR}_{\text{BPE}}$ при
$\Delta \mathrm{BPN} \leq 0{,}01$, что является более строгим
критерием.

\section{Аблационное исследование}

Аблационные эксперименты разделены на принципиальные и вспомогательные
по степени их влияния на архитектурные решения. В принципиальной
группе сравнивается жесткое DP с BCE-supervision (базовая
конфигурация) с альтернативными стратегиями обратного
распространения~--- log-semiring
relaxation~\cite{menschDifferentiableDynamicProgramming2018} и
Straight-Through Estimator~\cite{bengioEstimatingPropagatingGradients2013}.
Также в этой группе оценивается переход от глобального DP к окопному с
шириной окна $W \in \{256, 512, 1024\}$, замена линейного слияния на
МЛП с произведением Адамара, чувствительность к ширине токенизатора
$d_{\text{tok}} \in \{32, 64, 128, 256\}$ и к параметру сжатия $r \in
\{0{,}2, 0{,}3, 0{,}4\}$.

Вспомогательная группа покрывает второстепенные компоненты: отключение
контекстуализатора (CausalConv1D), отключение length embedding,
чувствительность к глубине рекурсии $L_{\max}$, замена per-bp декодера
на полноценное перекрестное внимание. Для каждой аблационной
конфигурации запуск выполняется с тем же зерном случайного генератора,
что и базовая конфигурация, что позволяет минимизировать влияние
случайной инициализации на сравнение.

\begin{table}[h]
\centering
\caption{Аблационное исследование~--- метрики предобучения и DART-Eval}
\label{tab:ablation}
\begin{tabular}{lccc}
\toprule
Конфигурация & PPL & BPN & DART-Eval (avg) \\
\midrule
Базовая (полная модель) & TBD & TBD & TBD \\
Замена жесткого DP на STE & TBD & TBD & TBD \\
Замена жесткого DP на log-semiring & TBD & TBD & TBD \\
Замена линейного слияния на МЛП с Hadamard & TBD & TBD & TBD \\
Замена per-bp декодера на cross-attention & TBD & TBD & TBD \\
Отключение CausalConv1D & TBD & TBD & TBD \\
Отключение length embedding & TBD & TBD & TBD \\
Окопный DP, $W = 256$ & TBD & TBD & TBD \\
Окопный DP, $W = 1024$ & TBD & TBD & TBD \\
\bottomrule
\end{tabular}
\end{table}

\section{Результаты на прикладных задачах}

\subsection{DART-Eval}

DART-Eval~\cite{patel2025dartevalcomprehensivednalanguage} является
основным бенчмарком для сравнительной оценки качества представлений.
Используется протокол с тремя уровнями адаптации: zero-shot
likelihood (без дообучения, с использованием логитов модели), линейное
зондирование с замороженным бэкбоном, и LoRA-дообучение с рангом 8.
В каждом из режимов оценка проводится отдельно для каждой из пяти
групп задач~--- обнаружение мотивов транскрипционных факторов,
предсказание связывания TF, доступности хроматина, регуляторной
активности, эффектов вариантов.

\begin{table}[h]
\centering
\caption{Результаты на DART-Eval. Метрика~--- AUROC, выше лучше}
\label{tab:dart}
\begin{tabular}{lccc}
\toprule
Задача & Char-T & BPE-T & Адаптивный \\
\midrule
Обнаружение мотивов TF (zero-shot) & TBD & TBD & TBD \\
Предсказание связывания TF (probed) & TBD & TBD & TBD \\
Доступность хроматина (probed) & TBD & TBD & TBD \\
Клеточно-специфичная регуляторная активность & TBD & TBD & TBD \\
Эффекты вариантов (caQTLs/dsQTLs) & TBD & TBD & TBD \\
\midrule
Среднее по бенчмарку & TBD & TBD & TBD \\
\bottomrule
\end{tabular}
\end{table}

\subsection{Дополнительные бенчмарки}

В дополнение к DART-Eval оценка проводится на трех наборах,
покрывающих различные классы прикладных задач. Genomic
Benchmarks~\cite{gresovaGenomicBenchmarksCollection2023} включает пять
задач бинарной классификации с длинами последовательностей 200--4000
п.н. (промоторы, энхансеры, регуляторные элементы), что позволяет
оценить качество представлений на относительно простых задачах с
вариативной длиной входа. Nucleotide Transformer downstream
tasks~\cite{dalla-torreNucleotideTransformerBuilding2025}~--- набор из
18 задач с длинами 200--600 п.н. (сайты сплайсинга, гистоновые
модификации, промоторы, энхансеры) и протоколом
LoRA-дообучения; данный бенчмарк рассматривается как кросс-проверка на
коротких последовательностях. BEND~\cite{tangBENDBenchmarkDNA2024}
ориентирован на биологически осмысленные задачи (gene finding,
enhancer annotation, chromatin accessibility, splice site, variant
effect) и использует протокол линейного зондирования поверх
замороженных эмбеддингов. Этот протокол является более чувствительным
к качеству самих представлений, нежели к выразительности дообученной
головы.

\section{Анализ выученной токенизации}

Помимо численных метрик предусмотрен биологический анализ выученной
токенизации, направленный на проверку того, согласуются ли границы
выделенных моделью спанов с известными функциональными границами
генома.

В первую очередь рассматривается распределение длин спанов $\ell_j$ на
тестовой хромосоме chr8. Сравнение проводится с распределением длин
BPE-токенов на той же выборке; при этом ожидается, что адаптивная
токенизация даст более широкое распределение с длинными хвостами,
поскольку модель не имеет фиксированного словаря и может выделять
более длинные регулярные участки в повторяющихся регионах. Если
выученное распределение оказывается практически дельта-функцией с
$\ell_j = \text{const}$~--- это означает коллапс адаптивности, что
является сигналом необходимости коррекции стратегии обучения.

Во-вторых, выученные границы спанов сопоставляются с биологическими
аннотациями. В качестве источников аннотаций были использованы три
базы данных. GENCODE~\cite{frankishGENCODE2019} предоставляет
референсную аннотацию генов и транскриптов, что позволяет оценить
совпадение границ спанов с экзон-интронными переходами и стартовыми/
терминирующими кодонами. ENCODE
cCREs~\cite{encodeExpandedRegistry2024} содержит реестр кандидатных
цис-регуляторных элементов, по которому можно оценить, попадают ли
границы спанов на края промоторов, энхансеров и сайтов CTCF.
JASPAR~2024~\cite{rauluseviciuteJASPAR2024} дает позиционные весовые
матрицы более чем 2000 транскрипционных факторов; для оценки
обогащения границ на их мотивах используется precision/recall
совпадения краев спанов с краями мотивов в пределах $\pm 1$~п.н.
относительно случайной базовой линии.

В-третьих, проверяется ранговая корреляция Спирмена между длинами
спанов $\ell_j$ и средней оценкой консервативности
phastCons~\cite{siepelEvolutionarilyConservedElements2005} на
покрываемом интервале $[a_j, b_j]$. При гипотезе о том, что в
функционально значимых регионах модель выделяет более короткие спаны
ради сохранения нуклеотидного разрешения, тогда как в межгенных
повторах допускает более длинные, ожидается отрицательная корреляция.
Стоит отметить, что данная гипотеза не является самоочевидной:
ничто в архитектуре явно не привязывает длину спана к функциональной
значимости, и возникновение подобной корреляции в результате
самосупервизионного обучения было бы свидетельством того, что задача
MLM сама по себе несет достаточный сигнал для биологически
осмысленной сегментации.

В дополнение к указанному анализу проводится стратификация выученных
длин по типам регионов (промотор, экзон, интрон, межгенный, повтор) и
сравнение средних длин внутри каждого класса; данная декомпозиция
позволяет выявить различия, которые могут быть скрыты на агрегированной
гистограмме.

\section{Обсуждение}

\subsection{Сопоставление с базовыми моделями}

При интерпретации результатов сопоставление с базовыми моделями
проводится по двум осям. Первая~--- метрики предобучения PPL и BPN при
сопоставимом числе параметров; данная ось отражает общее качество
модели как языковой модели нуклеотидов. Вторая ось~--- результаты на
прикладных задачах после дообучения; здесь существенным является не
абсолютная точность, а ее изменение относительно базовых моделей при
переходе к адаптивной токенизации. Раздельное рассмотрение этих осей
имеет значение по следующей причине: посимвольная HyenaDNA проигрывает
$k$-мерному Nucleotide Transformer~\cite{dalla-torreNucleotideTransformerBuilding2025}
по числу параметров, но превосходит его на ряде задач
зондирования~\cite{vishniakovGenomicFoundationlessModels2024}~--- то
есть качество представлений и качество дообученной головы могут
двигаться в разных направлениях. Адаптивная токенизация позиционируется
как промежуточный режим, в котором сжатие достигается без потери
посимвольной выразительности на функционально значимых границах.

\subsection{Стратегия обратного распространения}

Сравнение жесткого DP с BCE-supervision против STE и log-semiring
relaxation является принципиальным с точки зрения научного интереса.
Жесткий DP с имитационным сигналом обладает следующими
особенностями: (а) точное решение задачи выбора непересекающихся пар
без аппроксимаций; (б) отсутствие смещения, связанного с релаксацией
$\max$; (в) зависимость от качества инициализации функции оценки на
старте обучения, поскольку случайные оценки приводят к случайному
выбору пар, и BCE-сигнал воспроизводит эту случайность. Log-semiring
relaxation, напротив, дает плотный неэвристический градиент через
soft-маргиналы $\alpha_i$, что является существенным преимуществом на
ранних этапах, но платит за это смещением (forward $\neq$ backward в
пределе $\beta \to \infty$). STE является промежуточным вариантом,
простым в реализации, но дающим заведомо смещенный градиент. В случае,
если разница в финальных метриках укладывается в шум по зернам, выбор
между этими стратегиями является вопросом инженерной простоты, а не
качества; в противном случае~--- одна из стратегий имеет существенное
преимущество.

\subsection{Линейное слияние и компенсация глубиной}

Сравнение линейного слияния с двухслойным МЛП с произведением Адамара
проверяет гипотезу о компенсации упрощения отдельного слияния глубиной
рекурсии. Если ухудшение от линейного слияния на одноуровневой схеме
($L_{\max} = 1$, $L = 2048$) велико, тогда как на многоуровневой
($L_{\max} \geq 3$, $L \geq 8192$) разница исчезает~--- гипотеза
подтверждается. В обратном случае~--- линейное слияние действительно
является упрощением, не компенсируемым глубиной, и для качества
необходим переход к МЛП-варианту.

\subsection{Per-bp декодер против cross-attention}

Замена per-bp декодера на полноценное перекрестное внимание~--- одна
из наиболее затратных аблаций по числу FLOPs ($O(LNd)$ против $O(Ld^2)$
у per-bp). Если разница в метриках на DART-Eval оказывается
несущественной, выбор per-bp как базовой конфигурации оправдан с
вычислительной точки зрения. В противном случае~--- per-bp декодер
теряет существенную часть информации, и cross-attention становится
обязательным компонентом, что меняет общую вычислительную картину
конвейера.

\subsection{Ограничения}

Ряд аспектов остается за рамками текущей работы. Во-первых, оценка
проводится только на одном референсном геноме (CHM13), и перенос на
бактериальные пангеномы или геномы растений требует отдельного
исследования из-за иного транспозонного ландшафта. Во-вторых,
размерность $d_{\text{tok}} = 64$ выбрана эмпирически, и при ablation
на $d_{\text{tok}} \in \{32, 128, 256\}$ возможно обнаружится
оптимальная точка вне выбранного значения. В-третьих, вспомогательная
функция потерь $\mathcal L_{\text{aux}}$ в базовой конфигурации
имеет нулевой вес, и ее ненулевое значение само является переменной
аблационного исследования; результирующий выбор зависит от того,
действительно ли BCE-сигнал улучшает сходимость функции оценки или
является избыточным. В-четвертых, значение целевого числа токенов
$L_{\text{target}}$ задано конфигурацией, а не выбирается
автоматически; автоматический подбор требует отдельной оптимизационной
постановки.

\section{Выводы по главе}

В данной главе зафиксированы условия и протокол экспериментального
сравнения предлагаемого метода с пятью базовыми архитектурами на
четырех бенчмарках; зафиксированы метрики предобучения (PPL, BPN, TR,
CR) и прикладные критерии успешности (DART-Eval, Genomic Benchmarks,
NT downstream, BEND). Аблационное исследование структурировано по
девяти конфигурациям с разделением на принципиальные и
вспомогательные, что позволяет изолировать вклад каждого
архитектурного решения. Биологический анализ выученной токенизации
дополняет численные результаты тремя видами наблюдений: распределение
длин спанов, обогащение границ на пересечении с GENCODE, ENCODE
cCREs и JASPAR, ранговая корреляция длин со phastCons. Численные
значения, полученные при прохождении полного экспериментального
цикла, заполнят таблицы~\ref{tab:pretrain}, \ref{tab:ablation} и
\ref{tab:dart}, после чего основные выводы будут уточнены и
сопоставлены с гипотезами, сформулированными в данной главе.
